\subsection*{First results: quick grid, union deposits, \code{main}}

\begin{center}
\scriptsize
\setlength{\tabcolsep}{3.5pt}
\begin{tabular}{lcccccccccc}
\toprule
$\psi$ & $Y_0$ fit & $Y_0$ exact & $\sum Y$ & $\sum I$ & $n_{D,0}$ & spread$_0$ & $r_{D,0}$ & $\Delta s_F$ q4 & nfa q8 & $Y_0$ exact\\
 & \% & \% & 25q & 25q & \% & bp/yr & bp/yr & pp & \% $Y$ & $\psi{=}2$ rules\\
\midrule
1 & -0.047 & -0.060 & +1.64 & +8.24 & -5.09 & +94.8 & -47.8 & -0.342 & +0.701 & -0.069 \\
\textbf{2} & -0.064 & -0.076 & +1.63 & +8.49 & -5.20 & +100.1 & -50.6 & -0.445 & +0.756 & -0.076 \\
4 & -0.051 & -0.062 & +1.90 & +8.98 & -5.20 & +102.3 & -52.3 & -0.505 & +0.745 & -0.079 \\
8 & -0.047 & -0.056 & +2.03 & +9.19 & -5.22 & +103.7 & -53.2 & -0.537 & +0.744 & -0.081 \\
32 & -0.043 & -0.052 & +2.13 & +9.36 & -5.23 & +104.9 & -53.9 & -0.563 & +0.743 & -0.083 \\
\midrule
range / $|m(2)|$ & 33\% &  & 30\% & 13\% & 3\% & 10\% & 12\% & 50\% & 7\% &  \\
\bottomrule
\end{tabular}
\end{center}
{\footnotesize Headline shock $p^d=2\%$/qtr, differenced IRF from each solution's own rest point.
Coarse grid (25 points), a preview: at $\psi=2$ the full grid gives $-0.073\%$ against
$-0.064\%$ here. The last column re-clears the period map exactly at impact under each $\psi$ with the
$\psi=2$ continuation rules held fixed (the within-period effect alone). The range row is
$(\max m-\min m)/|m(\psi{=}2)|$ over the solved values. No solution at $\psi=0.5, 0.25$: the warm walk (1, 2 and 4 log-steps from $\psi=1$) and the cold ladder both fail.}

\begin{center}
\small
\begin{tabular}{lccccc}
\toprule
dial (values swept) & $\ln$ range & $Y_0$ range & $\sum Y$ range & $\Delta s_F$ range & nfa q8 range\\
\midrule
$\psi$ (1--32) & 3.47 & 33\% & 30\% & 50\% & 7\%\\
$\kappa_{\rm nfa}$ (0.0025--0.04) & 2.77 & 110\% & 56\% & 3\% & 9\%\\
\addlinespace
\multicolumn{6}{l}{per unit of $\ln$(dial): $Y_0$: $\psi$ 9.5\%, $\kappa_{\rm nfa}$ 40\%; \quad $\sum Y$: $\psi$ 8.7\%, $\kappa_{\rm nfa}$ 20\%}\\
\bottomrule
\end{tabular}
\end{center}
{\footnotesize Same grid, shock and reads as above; ranges as a share of the value at the
calibrated point. $\kappa_{\rm nfa}$: impact output $-0.085$, $-0.064$, $-0.014\%$ at
$0.0025$, $0.01$, $0.04$; cumulative output $+2.02$, $+1.63$, $+1.10$
(\code{results/STRESS/kappa\_nfa\_default\_s0.json}).}

\paragraph{Reading.}
\begin{enumerate}
\item \emph{Solvability bounds $\psi$ from below.} No solution exists at $\psi\le0.5$: the warm
walk from $\psi=1$ (in one, two and four log-steps) stalls at $\max|F|\approx3\cdot10^{-2}$, and
the cold ladder produces a non-finite iterate in the default-branch haircut homotopy. In that
branch the F bank wants to sell its whole D book; with the $10^{-4}$ floors in place of a
complementarity pair the Jacobian degenerates. This is the national branch's failure, reached
here through the portfolio cost. The interesting low-$\psi$ region cannot be explored without
porting the no-short-sale pairs to the union economy.
\item \emph{Under union funding, $\psi$ governs composition, not the current account.} The F
share moves by half its value across the sweep, while $D$'s net external position moves by 7\%,
because the deposit position $V$ absorbs the external adjustment whatever $\psi$ is (the
six-panel figure, \code{results/STRESS/psi\_bF\_D\_default\_s0.png}). This is the prediction of
Section~\ref{sec:diag}, job (c), read in reverse.
\item \emph{Output: a third over a 32-fold range, with a bump at the calibrated point.} A
stiffer friction gives a smaller impact fall and a larger later boom. Above $\psi=2$ both move
monotonically: the one-sided elasticity between 2 and 32 is $-0.15$ for $|Y_0|$ and $+0.10$
for cumulative output. The $\psi=1$ point breaks monotonicity on impact. That is \emph{not}
multiplicity: cold-solving $\psi=4$ reproduces the warm solution exactly ($-0.0514\%$), and
walking $4\to2$ warm reproduces the cold $\psi=2$ ($-0.0644\%$). On 25 points it may be an
interpolation effect at the KKT kink; the full-grid points $\{1,2,4\}$ decide.
\item \emph{The mechanism runs through expectations.} Holding the $\psi=2$ rules fixed and
changing $\psi$ only in the period map moves impact output the \emph{other} way ($-0.069\%$ at
$\psi=1$ to $-0.083\%$ at $\psi=32$): today a stiffer friction stops F banks buying at the dip,
so the D bank absorbs more. The fully re-solved effect has the opposite sign, so what $\psi$
does to the IRF is mostly what it does to expected future rebalancing, which a fixed-rule read
excludes.
\item \emph{The level anchor shows up as predicted.} The F share falls \emph{more} as $\psi$
rises ($-0.34$ to $-0.56$ pp at q4), the opposite of what a stiffer rebalancing friction should
do, because the friction holds the level while the debt stock grows. At $\psi=1$ the F share
rises on impact (F banks buy at the dip) before falling.
\item \emph{The funding closure dominates.} Per log unit, $\kappa_{\rm nfa}$ moves impact output
four times as much as $\psi$, and cumulative output twice as much. A stationarity device borrowed
from the small-open-economy literature, never calibrated here, is the most powerful dial of the
two. Its effect does not run through the loan rate (impact $r^{wc}$ is $+50$ bp/yr at
$\kappa=0.0025$ and $+48$ at $0.04$) but through the terms of trade and labour supply, as the
external margin tightens. The deposit-rate channel of the decomposition ($+0.104$ pp) is likewise
about five times the whole $\psi$ range of impact output (0.021 pp). For the \emph{union}
economy, then, the answer to Issue~2 is that $\psi$ drives the cross-border composition strongly
and the output path moderately, and Issue~1 is first-order. Under national funding the
prediction is the reverse; that sweep is the next run.
\end{enumerate}
